\SourceLeaf{294}%
\[
\begin{aligned}
t&=t^0(1-\tfrac{1}{6}\rho\rho rr-\tfrac{1}{3}\rho\rho t^0t^0)\\
\lambda&=\lambda^0(1-\tfrac{1}{6}\rho\rho s^0s^0-\tfrac{1}{3}\rho\rho t^0t^0)\\
\sigma&=\tfrac{1}{2}\rho vt^0(1-\tfrac{1}{12}\rho\rho rr-\tfrac{1}{4}\rho\rho s^0s^0-\tfrac{1}{4}\rho\rho t^0t^0)\\
\tau&=\tfrac{1}{2}\rho vs^0(1+\tfrac{5}{12}\rho\rho rr-\tfrac{1}{2}\rho\rho s^0s^0)
\end{aligned}
\]
accurate up to, but excluding, fifth-order quantities for \(t\) and \(\lambda\), and sixth-order quantities for \(\sigma\) and \(\tau\). It is even more convenient and just as accurate to use logarithms from the outset. If for brevity we denote the product of the quantity \(\tfrac{1}{12}\rho\rho\) and the modulus of Briggs's logarithms by \(\mu\), the formulas take the following form:
\[
\begin{aligned}
\log s&=\log s^0+4\mu rr-4\mu s^0s^0\\
\log t&=\log t^0-2\mu rr-4\mu t^0t^0\\
\log\lambda&=\log\lambda^0-2\mu s^0s^0-4\mu t^0t^0\\
\log\sigma&=\log\tfrac{1}{2}\rho vt^0-\mu rr-3\mu s^0s^0-3\mu t^0t^0\\
\log\tau&=\log\tfrac{1}{2}\rho vs^0+5\mu rr-6\mu s^0s^0
\end{aligned}
\]
\indent These five formulas, together with the preceding ones for \(s^0,t^0,\lambda^0\), form a fifth method of solution. Its distinctive feature is that approximate values of \(s,t,\lambda,\sigma,\tau\) are made accurate by small corrections to the logarithms, which are very easy to calculate. The constant logarithms used here are
\[
\begin{aligned}
\log\rho&=4{,}6855748668&&(-10)\\
\log\tfrac{1}{2}\rho&=4{,}3845448712&&(-10)\\
\log\mu&=7{,}9297527989&&(-20)
\end{aligned}
\]
or, if the corrections are to appear directly in units of the seventh decimal place,
\[
\log\mu=4{,}9297527989\qquad(-10)
\]
Only the first digits of these logarithms are used here, however.
\Article{17}
The relations between the quantities \(r,S,S',T,T',\lambda\) can be expressed much more simply, however, if we start from the mean of the two latitudes
